% Bagian: second-order-logic
% Bab: sol-and-sets
% Subbagian: comparing-sets

\documentclass[../../../include/open-logic-section]{subfiles}

\begin{document}

\olfileid[id]{sol}{set}{cmp}

\olsection{Membandingkan Himpunan}

\begin{prop}
!!^{formula} $\lforall[x][(X(x) \lif Y(x))]$ mendefinisikan relasi himpunan
bagian, yakni $\Sat{M}{\lforall[x][(X(x) \lif Y(x))]}[s]$ jika dan hanya jika
$s(X) \subseteq s(Y)$.
\end{prop}

\begin{prop}
!!^{formula} $\lforall[x][(X(x) \liff Y(x))]$ mendefinisikan relasi identitas
pada himpunan, yakni $\Sat{M}{\lforall[x][(X(x) \liff Y(x))]}[s]$ jika dan
hanya jika $s(X) = s(Y)$.
\end{prop}

\begin{prop}
!!^{formula} $\lexists[x][X(x)]$ mendefinisikan sifat tidak kosong, yakni
$\Sat{M}{\lexists[x][X(x)]}[s]$ jika dan hanya jika $s(X) \neq
\emptyset$.
\end{prop}

Suatu himpunan~$X$ tidak lebih besar daripada himpunan~$Y$,
$\cardle{X}{Y}$, jika dan hanya jika terdapat fungsi !!a{injective}
$f\colon X \to Y$. Karena kita dapat mengungkapkan bahwa suatu fungsi
bersifat injektif, dan juga bahwa nilai-nilainya untuk argumen dalam~$X$
berada dalam~$Y$, kita juga dapat mendefinisikan relasi tidak lebih besar
daripada di antara himpunan-himpunan bagian domain.

\begin{prop}
Formula
\[
\lexists[u][(\lforall[x][(X(x) \lif Y(u(x)))] \land
  \lforall[x][\lforall[y][((X(x) \land X(y) \land
    \eq[u(x)][u(y)]) \lif \eq[x][y])]])]
\]
mendefinisikan relasi tidak lebih besar daripada.
\end{prop}

Dua himpunan berukuran sama, atau ``berkardinalitas sama,''
$\cardeq{X}{Y}$, jika dan hanya jika terdapat fungsi
!!a{bijective}~$f\colon X \to Y$.

\begin{prop}
Formula
\begin{multline*}
  \lexists[u][(\lforall[x][(X(x) \lif Y(u(x)))] \land {}\\
    \lforall[x][\lforall[y][((X(x) \land X(y) \land
      \eq[u(x)][u(y)]) \lif \eq[x][y])]]
  \land {} \\\lforall[y][(Y(y) \lif \lexists[x][(X(x)
      \land \eq[y][u(x)])])])]
\end{multline*}
mendefinisikan relasi berkardinalitas sama.
\end{prop}

Kita akan menyingkat !!{formula}s ini, secara berturut-turut, sebagai
$X \subseteq Y$, $X = Y$, $X \neq \emptyset$, $\cardle{X}{Y}$, dan
$\cardeq{X}{Y}$. (Ini mungkin sedikit membingungkan karena kita menggunakan
notasi yang sama ketika berbicara secara informal tentang himpunan $X$ dan
$Y$---tetapi di sini notasi tersebut merupakan singkatan bagi !!{formula}s
dalam logika orde kedua yang melibatkan variabel relasi satu-tempat $X$
dan~$Y$.)

\begin{prop}
!!^{sentence} $\lforall[X][\lforall[Y][((\cardle{X}{Y} \land
    \cardle{Y}{X}) \lif \cardeq{X}{Y})]]$ valid.
\end{prop}

\begin{proof}
!!^{sentence} tersebut dipenuhi dalam !!a{structure}~$\Struct{M}$ jika untuk
sebarang himpunan bagian $X \subseteq \Domain{M}$ dan
$Y \subseteq \Domain{M}$, jika $\cardle{X}{Y}$ dan $\cardle{Y}{X}$, maka
$\cardeq{X}{Y}$. Namun, hal ini berlaku bagi \emph{sebarang} himpunan $X$
dan $Y$---inilah Teorema Schr\"oder--Bernstein.
\end{proof}


\end{document}
